% notation.tex -- T1 catalogue, gate G0.
% One table for the paper's three naming systems, one name per object, units everywhere.
% Compiles standalone (latexmk -pdf notation.tex) and is safe to \input into the paper body.
\documentclass[11pt]{article}
\usepackage[margin=2.2cm]{geometry}
\usepackage{longtable,array,amsmath,amssymb}
\usepackage[T1]{fontenc}
\newcolumntype{L}[1]{>{\raggedright\arraybackslash}p{#1}}
\newcommand{\C}{\mathbb{C}}
\newcommand{\R}{\mathbb{R}}
\begin{document}

\section*{Notation: one name per object}

These are binding rules for the final file (T8) and for the checker's record format (T2).

\begin{enumerate}
\item \textbf{One name per object.} Where the paper uses one letter for two things, this file picks
  one meaning and renames the other. Table~2 lists the actual collisions.
\item \textbf{Units are part of the statement.} Every numerical constant carries its unit:
  \emph{cycles} (phase divided by $2\pi$, the unit in which radii such as $0.032$, $0.125$, $.016$,
  $.011$, $.0038$ and $.042$ are stated) or \emph{radians} (the unit in which $\theta$, $h=\pi/N_u$,
  $s$ as a seed displacement, $d_j$ as an angular displacement and arc are stated), or
  \emph{dimensionless}. The conversions $R_s = 2\pi r_s$ and the $/(2\pi)$ and $/\pi$ in
  Lemma~4.10 are the only places where the two mix, and both must be written with their factor.
\item \textbf{The seven Root constants are always written in Definition~4.1 form.}
  Sections 5 and 6 rename them; those renamings are translation tables, not notation.
\end{enumerate}

\subsection*{Table 1: the seven Root constants, three naming systems}

Definition~4.1 introduces them as $C,n_0,\mathrm{id},\mathrm{ort},\mathrm{hor},\mathrm{mnorm},\mathrm{froot}$;
\S5.1 renames five of them; \S6.1 renames four of them again; \S7.1 stores the first set in every record.

\begin{longtable}{@{}L{1.05cm}L{1.15cm}L{1.2cm}L{4.9cm}L{1.9cm}L{0.9cm}@{}}
\hline
Def.~4.1 & \S5.1 & \S6.1 & what it is, in words & type / unit & first \\
\hline
\endhead
$C$ & $\gamma_P$ & $\Gamma$ &
Coefficient of $s^2$ in the normal polynomial $n_P(s)$: how fast the bound on the
normal (off-tangent) part of a displacement grows quadratically with its total size. Certified as
$\pi$ times a remainder constant $b$ (Lemma~4.3, 4.4). &
nonneg.\ real, dim.\ free & \S4.1 \\
$n_0$ & $n_0$ & $n_0$ &
Constant term of $n_P$: the normal-displacement bound in the limit of a zero-size displacement.
Certificate requires $n_0 \ge \mathrm{mnorm}\cdot\mathrm{froot}$ (4.15); \S6.1 admits $n_0<10^{-6}$. &
nonneg.\ real, cycles & \S4.1 \\
$\mathrm{id}$ & $\alpha_P$ & $\iota$ &
Coefficient of $s$ in $n_P$: the linear growth rate of the normal bound. Comes from the certified
deviation identity (4.15); \S6.1 admits $\mathrm{id}<10^{-4}$. &
nonneg.\ real, dim.\ free & \S4.1 \\
$\mathrm{ort}$ & $o_P$ & $\omega$ &
Orthogonality defect of the stored tangent frame: $\|T^*T-I\|$. Bounds the loss when a displacement is
reconstructed from its tangent coordinates. \S6.1 admits $0\le\omega<10^{-5}$, a site of \S6.1's
unclosed hypothesis (see report item I3). &
nonneg.\ real, dim.\ free & \S4.1 \\
$\mathrm{hor}$ & $\mathrm{hor}_P$ & --- &
Horizontality defect of the frame: $\|(I-\Pi)T\|$. Not renamed in \S5.1, unused in \S6.1.
The only Root constant that keeps one name throughout. &
nonneg.\ real, dim.\ free & \S4.1 \\
$\mathrm{mnorm}$ & $\mu_P$ & --- &
Operator norm bound for the auxiliary matrix $M$ in (4.15): $\|M\|\le\mathrm{mnorm}$, used with
$\mathrm{mnorm}\cdot\mathrm{froot}\le n_0$. &
nonneg.\ real, dim.\ free & \S4.1 \\
$\mathrm{froot}$ & $\varepsilon_P$ & $f_0$ &
Bound on the residual at the stored centre: $\|F(P)\|\le\mathrm{froot}$. This is the
``distance of the centre from the Hadamard locus''; \S6.1 admits $f_0<10^{-6}$. &
nonneg.\ real, dim.\ free & \S4.1 \\
\hline
\end{longtable}

\noindent\emph{Rule.} After this table the paper writes these seven always as
$C,n_0,\mathrm{id},\mathrm{ort},\mathrm{hor},\mathrm{mnorm},\mathrm{froot}$, with the centre as a
subscript when two Roots are in play: $C_P$, $n_{0,P}$, $\mathrm{id}_P$, $\mathrm{ort}_P$,
$\mathrm{hor}_P$, $\mathrm{mnorm}_P$, $\mathrm{froot}_P$. The second Root in \S5 is $Q$, with its
tangent matrix $S$ (not $T$).

\subsection*{Table 2: letters that carry two or three objects in the paper}

\begin{longtable}{@{}L{1.1cm}L{3.4cm}L{5.6cm}L{2.2cm}@{}}
\hline
letter & the paper's uses & the fix & where \\
\hline
\endhead
$D$ &
(i) the $4\times4$ comparison block $D=D_{\mathrm{span}}+D_{\mathrm{ker}}$ of \S3.4;
(ii) a scalar distance upper bound in Prop.~5.2 ($D\ge R+\|d_q\|$);
(iii) $\hat D(t)$, the computed Stage-B function of \S5.2. &
Write $D^{\mathrm{cmp}}$ for (i) (\emph{comparison block}), $D^{\mathrm{ub}}$ for (ii), and keep
$\hat D(t)$ for (iii) with its argument always displayed. The comparison block is not the lower-right
block of a Hadamard matrix at the centre; say so once. &
\S3.4, \S5.1, \S5.2 \\
$n_P$ &
(i) the normal polynomial $n_P(s)=n_0+\mathrm{id}\,s+Cs^2$;
(ii) the normal displacement bound $n=n_0+\iota r+\Gamma r^2$ of (6.2). &
$n_P$ is the function, $n_P(s)$ its value. In \S6 write $n_P(r)$, never a bare $n$. The word
\emph{normal} always means ``the part of $\delta$ orthogonal to the stored tangent frame'', so
$n_P$ is a \emph{bound}, not a coordinate. &
\S4.1, \S6.1 \\
$m$ &
(i) the pair of lower singular bounds $m_B=\sigma_{\min}(B_0)$, $m_C=\sigma_{\min}(C_0)$ and the
shorthand $m=\min(m_B,m_C)>0$;
(ii) $m_j=|D_j|$, the moduli of the sixteen comparison-block entries (Lemma~4.8);
(iii) $m\in\mathbb{Z}^6$, the integer lift vector in Stage-C clustering (Lemma~6.4). &
Keep $m_B,m_C,m$ for (i) but always with the subscript or the defining display;
write $\dot m_j := |D_j|$ (or spell it) for (ii); write $\mathbf{m}\in\mathbb{Z}^6$ for (iii). &
\S3.1, \S4.4, \S6.3 \\
$\varepsilon$ &
(i) $\varepsilon$, the stored-centre correction in cycles, $\|\varepsilon\|\le2\cdot10^{-6}$ (4.20);
(ii) $\varepsilon_P=\mathrm{froot}_P$ (\S5.1);
(iii) $\varepsilon_{\exp}=6\cdot10^{-14}$, the unit-exponential error (A.10). &
Keep $\varepsilon$ for (i) only (it is a vector in $\R^{36}$); write $\mathrm{froot}_P$ for (ii);
write $\varepsilon_{\exp}$ with the subscript for (iii). &
\S4.4, \S5.1, \S A.2 \\
$N$ &
(i) the seed grid resolution $N=112$, $w=1/N$;
(ii) an integer $6\times6$ lift matrix in (4.17), Lemma~4.15. &
Write $N_{\mathrm{grid}}$ (or just $112$) for (i); write $N^{\mathrm{lift}}$ or $\mathbf{N}$ for (ii),
and always display the integrality hypothesis next to it. &
\S2.2, \S4.3 \\
$k$ &
(i) the number of stored tangent columns, $0\le k\le10$;
(ii) $k=N(u-u')\in\mathbb{Z}^9$, the integer grid difference in Lemma~2.9. &
$k$ stays (i); write $\kappa$ for (ii). &
\S4.1, \S2.4 \\
$s$ &
(i) the nine-component seed displacement $s\in\R^9$ (\S3.1) and its norm $s=\|\delta\|$ in Def.~4.1;
(ii) the radius $s$ of a Stage-C vector ball $B_{\mathrm{vec}}(v,s)$. &
In Definition~4.1 use $s=\|\delta\|$ exactly as printed, but in \S3 write $s$ only for the vector
(it is the symbol the code stores); in \S6.3 write $s$ for a ball radius and nothing else. Never
let a bare $s$ denote both a vector and a scalar in one paragraph. &
\S3.1, \S4.1, \S6.3 \\
$d$ &
(i) $d_q=\Pi(D_0(Q)-D_0(P)-N)$, a displacement vector;
(ii) $d_j$, the sixteen angular displacements in radians (Lemma~4.8);
(iii) the scalar $d=\|\delta-d_q\|$ of Prop.~5.2 and $\hat D$. &
$d_q$ keeps its subscript; $d_j$ stays; write $d$ only for the scalar, and state
$d=\|\delta-d_q\|$ each time. &
\S4.3--\S5.2 \\
$z$, $\zeta$ &
$z=\|Z\|_F$ (comparison residual bound, \S3.4); $z=x_u-x_v$ (edge vector, \S6.4);
$\zeta=T^*x$ (Lemma~5.1). &
$z$ stays a nonnegative scalar bound in \S3--\S4; in \S6.4 write $z_e$ for the edge vector;
$\zeta$ is a vector and is always defined where used. &
\S3.4, \S5.1, \S6.4 \\
$S$ &
$S$ is the Q-Root tangent matrix (\S5.1), the block $S$ of a vector sum (\S6.2), and the $30\times216$
tensor map of Lemma~4.4. &
Write $T_Q$ for the Q-Root tangent matrix (and then $T$ is $T_P$), $S_i=\sum_j q_{ji}$ for the
five-vector sums, and $S^{\mathrm{ten}}$ for the tensor map of Lemma~4.4. &
\S4.2, \S5.1, \S6.2 \\
$r$ &
radii (cycles) and the scalar remainder $r(t)=(e(t)-1-2\pi it)/(2\pi i)$ of Lemma~4.3. &
$r$ for radii only; $r_{\mathrm{rem}}(t)$ for the scalar remainder. Radii: $R$ for a seed radius in
radians, $r_s=R/(2\pi)$ for the same radius in cycles, $r$ for a Record radius in cycles. &
\S2.4, \S4.1, \S4.3 \\
\hline
\end{longtable}

\subsection*{Table 3: the words}

The paper uses six geometric words and two capitalised data words. Fix them once.

\begin{longtable}{@{}L{1.9cm}L{9.6cm}L{2.0cm}@{}}
\hline
word & fixed meaning & first \\
\hline
\endhead
\textbf{bin} & a closed axis-aligned box in phase coordinates, of width $1/N$ in \S2 ($N=112$) or
  $1/M$ in \S6; several coordinates at once form a \emph{bin tuple}. Bins are the objects the
  enumeration visits. & \S2.2, \S6.2 \\
\textbf{cell} & a closed axis-aligned box in the four angle coordinates $(\alpha,t,\beta,\gamma)$ of
  $U(2)$ (\S3.3), of width $1/N_u$ per coordinate, with midpoint $U$. Cells are what Stage A
  subdivides and either excludes or covers. Use \emph{cell} only for this; never for a phase box. & \S3.3 \\
\textbf{box} & a closed axis-aligned box in the $k$ tangent coordinates (\S5.2), with centre $u$ and
  halfwidth $h$; also called a \emph{tangent box}. A box is marked only if the whole box passes the
  test (Lemma~5.5); a successful subbox does not license marking its base box. & \S5.2 \\
\textbf{ball} & a closed Euclidean ball, in one of two spaces: a \emph{phase ball} is
  $\{\delta\in V:\|\delta\|\le r\}$ around a phase centre (radius in cycles), a \emph{vector ball}
  $B_{\mathrm{vec}}(v,s)$ is a set of phase vectors up to global phase (radius in cycles). Say which. & \S2.4, \S4.7, \S6.3 \\
\textbf{sphere} & the word is used in the paper, and in the source, for exactly the closed phase ball
  of a Record. It is not a topological sphere and it includes its boundary. \emph{Replace every
  occurrence by ``ball''}; do not use ``sphere'' at all. & \S4.7, Def.~4.14 \\
\textbf{Root} & the eight-tuple $(P,T,(C,n_0,\mathrm{id},\mathrm{ort},\mathrm{hor},\mathrm{mnorm},\mathrm{froot}))$
  of Definition~4.1: a centre plus a tangent frame plus seven bounds. A Root carries no radius. & \S4.1 \\
\textbf{Record} & a certified Root together with a radius $r$, i.e.\ a ball with a certificate. A
  Record is what the atlas stores and what Stage~C judges. & \S4.7, Def.~4.14 \\
\textbf{atlas} & the finite set of Records accumulated by Stage~A and refined by Stage~B. & \S4.7 \\
\hline
\end{longtable}

\subsection*{Table 4: other symbols used throughout}

\begin{longtable}{@{}L{1.6cm}L{5.6cm}L{2.4cm}L{1.0cm}@{}}
\hline
symbol & what it is, in words & type / unit & first \\
\hline
\endhead
$\Pi$ & double centring: the orthogonal projection onto matrices with zero row and column sums (the
  \emph{horizontal} subspace $V$). & lin.\ map $\R^{36}\to V$ & \S2.1 \\
$F(P)$ & the residual: $F_{ij}(P)=\frac{1}{2\pi i}\sum_g e(P_{ig}-P_{jg})$ for $i>j$, a 30-vector.
  $F(P)=0$ iff $e(P)$ is Hadamard. & $\R^{36}\to\R^{30}$ & \S4.2 \\
$D_0(P)$ & dephasing: $D_0(P)_{ij}=P_{ij}-P_{i0}-P_{0j}+P_{00}$, first row and column zero. & $\R^{36}\to\R^{36}$ & \S2.1 \\
$L_P$ & first variation of $F$ at $P$. & lin.\ map & \S4.2 \\
$R_P(x)$ & quadratic remainder, $F(P+x)=F(P)+L_Px+R_P(x)$, $|R_P(x)_{ij}|\le\pi a_{ij}(x)$. & $\R^{30}$ & \S4.2 \\
$T$ & stored tangent frame of a Root, $36\times k$, $0\le k\le10$; columns are numerical directions
  only, neither exactly horizontal nor exactly orthogonal. & matrix & \S4.1 \\
$\eta,\nu$ & tangent and normal parts of a displacement: $\eta=T^*\delta$, $\nu=\delta-T\eta$
  (orthogonal decomposition by $TT^*$). & vectors & \S6.1 \\
$r_s$, $R_s$ & the seed radius in cycles ($r_s=0.0320001$) and in radians ($R_s=2\pi r_s$). & scalar & \S2.4, \S3.1 \\
$r_H$ & inflated tangent radius $\sqrt{1+\omega}\,r$ (6.2). & cycles & \S6.1 \\
$X_M,L_M,I_M$ & Stage-C bin displacement bounds: $\sqrt{29/6}\,a$, $5a$, $\tfrac32 a$ with $a=1/(2M)$. & cycles & \S6.2 \\
$a$ & half bin width $1/(2M)$ at resolution $M$. & cycles & \S6.2 \\
$K$ & operator bound for $e(P)/\sqrt6$: $K=1.00001$ suffices when $\mathrm{froot}<10^{-6}$. & scalar & \S6.1 \\
$q_d$ & normalised slack $\ge \mathrm{slack}/m_i$, required $<3.8$ (Lemma~4.8). & scalar & \S4.4 \\
arc, weighted & angular-displacement bounds of Lemma~4.8. & radians & \S4.4 \\
$E_r$ & linearisation remainder bound, $\|Z-J(s,\theta)\|_F\le E_r$ (3.24). & scalar & \S3.5 \\
$e_{\mathrm{off}}$ & the seed-product part of $E_r$ (3.25). & scalar & \S3.5 \\
$z$ & $\|Z\|_F$ bound (3.19). & scalar & \S3.4 \\
$\mathcal{O},\mathcal{U}$ & the orthogonality and unbiasedness graphs on retained ball labels; an edge
  is retained when the test cannot rule the relation out. & graphs & \S6.6 \\
$A,B$ (sets) & the two disjoint six-cliques of configuration (6.24); distinct from the $2\times4$
  seed blocks $A$--$D$. & sets of labels & \S6.6 \\
$E_{\mathcal{O}},E_{\mathcal{U}}$ & edge sets of the two graphs. & sets & \S6.6 \\
\hline
\end{longtable}

\subsection*{Macros}

\begin{verbatim}
% units
\newcommand{\cyc}{\,\text{cycles}}
\newcommand{\rad}{\,\text{rad}}
% the seven Root constants, Definition 4.1 names only
\newcommand{\rtC}{C}\newcommand{\rtnz}{n_0}\newcommand{\rtid}{\mathrm{id}}
\newcommand{\rtort}{\mathrm{ort}}\newcommand{\rthor}{\mathrm{hor}}
\newcommand{\rtmn}{\mathrm{mnorm}}\newcommand{\rtfr}{\mathrm{froot}}
% the renamings, for the translation table and nowhere else
\newcommand{\gP}{\gamma_P}\newcommand{\aP}{\alpha_P}\newcommand{\oP}{o_P}
\newcommand{\mP}{\mu_P}\newcommand{\eP}{\varepsilon_P}
\newcommand{\Gam}{\Gamma}\newcommand{\iot}{\iota}\newcommand{\omg}{\omega}
% resolved collisions
\newcommand{\Dcmp}{D^{\mathrm{cmp}}}     % comparison block
\newcommand{\Dub}{D^{\mathrm{ub}}}       % scalar distance bound, Prop 5.2
\newcommand{\Dhat}{\hat D}               % computed Stage-B function
\newcommand{\nP}{n_P}                    % normal polynomial
\newcommand{\mB}{m_B}\newcommand{\mC}{m_C}
\newcommand{\mj}{\dot m_j}               % |D_j|
\newcommand{\mlift}{\mathbf{m}}          % integer lift in Z^6
\newcommand{\Ngrid}{N_{\mathrm{grid}}}   % 112
\newcommand{\Nlift}{\mathbf{N}}          % integer lift matrix
\newcommand{\kappa}{\kappa}              % N(u-u') in Z^9
\newcommand{\eps}{\varepsilon}           % centre correction, cycles
\newcommand{\eexp}{\varepsilon_{\exp}}   % 6e-14
\newcommand{\TQ}{T_Q}                    % Q-Root tangent matrix
\newcommand{\rrem}{r_{\mathrm{rem}}}     % scalar remainder of Lemma 4.3
\newcommand{\zv}{z_e}                    % edge vector in 6.4
\newcommand{\ball}{\mathcal{B}}\newcommand{\vball}{B_{\mathrm{vec}}}
\end{verbatim}

\end{document}
